% ===================================================================
%  図表第 14 号 — 街路。— 商人たち。
%
%  Ceci n'est PAS une transcription : c'est une traduction, faite en
%  2026, en japonais d'aujourd'hui. Voir texto/ja/10-zuhyo-01.tex
%  pour la consigne, qui vaut pour les seize fichiers.
%
%  LE GENITIF JAPONAIS RENVERSE LES RENVOIS. « le jet d'eau (49) de la
%  fontaine (48) » se dit 噴水の水柱 — 48 avant 49 — et le releve
%  sortait faux. Sept fois dans ce seul tableau, on a donc coupe le
%  genitif par une incise : 水柱 (49) ——噴水 (48) の水柱—— . L'ordre
%  reste celui de l'ido, et la phrase reste japonaise.
% ===================================================================

%%K t14-apar-1 apar
\VUsaut{4.0mm}
\VUtitre{15.0pt}{60}{図表第 14 号}
\VUsaut{3.0mm}

%%K t14-tit-1 sub
\VUpk{11.4pt}{40}{街路。--- 商人たち。}
\VUsaut{3.0mm}

%%K t14-01-1 p
1. --- 田舎に住んでいる私の友だちのジャンが、うちに三日泊まりにきた。
\VUgras{それまで}大きな町を見る機会がなかったので、私たちは一日じゅう歩きまわり、
彼の知らないものを私が説明してやった。

%%K t14-02-1 p
2. --- まず\VUgras{天文台} \textsuperscript{(1)} を見に行った。彼はたいそう
面白がった。帰りは\VUgras{街路} \textsuperscript{(3)} ——\VUgras{トルコ風呂}
\textsuperscript{(2)} のある街路である——を通り、\VUgras{葬式} \textsuperscript{(4)} に
出会った。\VUgras{葬列}の先頭には\VUgras{聖職者}が進み、その後ろに
\VUgras{霊柩車}が、\VUgras{柩}を載せて続いた。大勢の友人が\VUgras{喪服の}家族に
付き添っていた。

%%K t14-02-2 p
一行が通りすぎてから、私たちは大きな\VUgras{ビヤホール} \textsuperscript{(5)} の前で
街路を横切った。大勢の\VUgras{客}がテラスで\VUgras{ビール}を飲んでいて、頭の上を
ガラス張りの\VUgras{日よけ} \textsuperscript{(6)} がおおっていた。

%%K t14-03-1 p
3. --- ジャンは\VUgras{紋章}に目を丸くした。\VUgras{酒屋} \textsuperscript{(7)} の店と
\VUgras{楽譜屋} \textsuperscript{(10)} の店のあいだに掛かっていたのである。彼はそれが何を
意味するのかと尋ねた。私はイギリスの\VUgras{領事館} \textsuperscript{(8)} を示すものだと
答え、バルコニーに立っている\VUgras{領事} \textsuperscript{(9)} を指してみせた。

%%K t14-tit-2 sub
\VUpk{10.2pt}{40}{飾り窓めぐり。}
\VUsaut{3.0mm}

%%K t14-04-1 p
4. --- もちろん私たちは\VUgras{楽器}、\VUgras{ハーモニウム}、\VUgras{オルガン}、
\VUgras{ピアノ}を並べた店先で足をとめた。\VUgras{楽譜}やピアノ用の\VUgras{曲集}の
色刷りの表紙を眺め、看板がわりの金箔を置いた木の\VUgras{竪琴}に見とれた。

%%K t14-05-1 p
5. --- \VUgras{清掃人} \textsuperscript{(11)} と\VUgras{街路清掃車}
\textsuperscript{(12)} の立てる土ぼこりのために、私たちは\VUgras{家具屋}
\textsuperscript{(13)} の立派な\VUgras{一そろいの家具}の前も、\VUgras{金物屋}
\textsuperscript{(14)} の並べた\VUgras{ストーブ}や\VUgras{ランプ}の前も、足早に通り
すぎるほかなかった。

%%K t14-06-1 p
6. --- しかもこの場所では、歩道をおりなければならなかった。\VUgras{運送車}
\textsuperscript{(16)} からおろして、借りたばかりの\VUgras{店} \textsuperscript{(15)} に
運びこむ荷物で、歩道がふさがっていたからである。

%%K t14-06-2 p
そのために\VUgras{馬車屋} \textsuperscript{(17)} の立派な車は見そこねたが、
\VUgras{荷造り人} \textsuperscript{(19)} が隣の店——\VUgras{自転車と自動車の店}
\textsuperscript{(18)} ——のために木箱を打っているのは、立ちどまって見ることができた。
そのあと少し遅くなったので、私たちは家に帰った。

%%K t14-tit-3 sub
\VUpk{10.2pt}{40}{町を歩く。}
\VUsaut{3.0mm}

%%K t14-07-1 p
7. --- 翌日、母がジャンの写真を撮ってもらってはどうかと言い出し、私に
\VUgras{写真屋} \textsuperscript{(20)} まで付き添うようにと言いつけた。昼食のあと私たちは
その職人の\VUgras{仕事場}に入り、そこでずいぶん長く待たされた。退屈しのぎに、壁に
掛かった\VUgras{肖像写真} \textsuperscript{(22)} を眺めた。

%%K t14-07-2 p
順番がくると仕事は長くはかからず、友だちが\VUgras{カメラ} \textsuperscript{(21)} の前で
ポーズを取り終えるとすぐ、私たちは下におりた。

%%K t14-08-1 p
8. --- 二階では、\VUgras{理髪師} \textsuperscript{(23)} の店の開いた戸口から、その
若い者が客の髪を刈っているのが見えた。

%%K t14-08-2 p
街路に出ると、私たちは\VUgras{たばこ屋} \textsuperscript{(24)} の前を通った。ジャンは赤い
\VUgras{提灯} \textsuperscript{(25)} と、\VUgras{看板} \textsuperscript{(26)} がわりの
\VUgras{大きなパイプ}とがあまり面白かったので、\VUgras{嗅ぎたばこを吸いながら}
\textsuperscript{(27)} 話しこんでいた二人の老紳士にぶつかってしまった。

%%K t14-tit-4 sub
\VUpk{10.2pt}{40}{菓子屋。}
\VUsaut{3.0mm}

%%K t14-09-1 p
9. --- \VUgras{両替屋} \textsuperscript{(29)} の\VUgras{飾り窓}に並べた諸国の
\VUgras{紙幣}や\VUgras{貨幣}にしばらく目を奪われたが、ちょうど間食の時刻だったので、
それ以上足をとめずに\VUgras{菓子屋} \textsuperscript{(31)} に入った。

%%K t14-10-1 p
10. --- 店にあるおいしそうな\VUgras{品}——\VUgras{ケーキ、蜂蜜菓子、ビスケット}、
それにいろいろの\VUgras{ボンボン}——を見て、私たちはなかなか選べなかった。とうとう
ジャンは\VUgras{クリーム菓子}を選び、私は小さな\VUgras{さくらんぼのタルト}だけに
した。甘いものは\VUgras{歯にさわる}し、\VUgras{歯医者} \textsuperscript{(30)} にたびたび
通うのはまっぴらだからである。

%%K t14-tit-5 sub
\VUpk{10.2pt}{40}{タイプライター。}
\VUsaut{3.0mm}

%%K t14-11-1 p
11. --- 話には聞いていても見たことのない\VUgras{タイプライター}を友だちに見せる
ため、私たちは\VUgras{事務室} \textsuperscript{(32)} に入った。\VUgras{従兄}
\textsuperscript{(34)} の事務室である。行ってみると、彼は\VUgras{秘書}
\textsuperscript{(35)} に手紙を口述しているところだった。秘書は\VUgras{速記者}でもある。
一通の手紙が終わるか終わらないうちに、\VUgras{タイピスト} \textsuperscript{(33)} がそれを
機械で清書した。親戚の仕事の邪魔をしたくなかったので、ほんの数分いただけだった。

%%K t14-12-1 p
12. --- 従兄の事務所の下には大きな\VUgras{時計宝石商} \textsuperscript{(37)} が住んで
いて、この人は銀細工師でもある。その立派な店には\VUgras{時計、置き時計、懐中時計}、
\VUgras{鎖}、それに高価な\VUgras{装身具}——たとえば\VUgras{真珠の首飾り}、金の
\VUgras{腕輪}、\VUgras{耳輪}、\VUgras{宝石}や\VUgras{ダイヤモンド}をはめた
\VUgras{指輪}——が並んでいる。飾り窓の一つはすっかり\VUgras{金銀細工}にあてられて
いて、銀の\VUgras{食器一式}が数多く納まっている。

%%K t14-13-1 p
13. --- 街角を曲がるとき、家の\VUgras{番地}のそばの\VUgras{標札}
\textsuperscript{(36)} が付け替えられているのに気づいた。それを声に出して言おうとした
とたん、軍楽の陽気な音が聞こえてきた。私たちは行列のほうへ駆け出したが、それが
\VUgras{帽子屋} \textsuperscript{(39)} と\VUgras{手袋屋} \textsuperscript{(40)} の店の前を
通るのだということも、\VUgras{眼鏡屋} \textsuperscript{(38)} の飾り窓——
\VUgras{鼻眼鏡、眼鏡}、\VUgras{オペラグラス}がぎっしり並んでいる——の前にいれば同じ
くらいよく見えるのだということも、考えなかった。

%%K t14-tit-6 sub
\VUpk{10.2pt}{40}{分列行進。}
\VUsaut{3.0mm}

%%K t14-14-1 p
14. --- \VUgras{連隊} \textsuperscript{(41)} の先頭を\VUgras{鼓手長}
\textsuperscript{(42)} が進み、そのあとに\VUgras{鼓手} \textsuperscript{(43)}、
\VUgras{らっぱ手} \textsuperscript{(44)}、そして\VUgras{楽隊}全体が続いた。
\VUgras{将軍} \textsuperscript{(45)} は副官とともに広場にいて、\VUgras{水柱}
\textsuperscript{(49)} ——\VUgras{噴水} \textsuperscript{(48)} の水柱である——のそばに
立ちどまり、\VUgras{分列}を見ていた。

%%K t14-14-2 p
\VUgras{参謀将校} \textsuperscript{(46)}、\VUgras{大佐}、\VUgras{中佐}は剣を挙げて将軍に
敬礼した。\VUgras{少佐}たちはそれぞれの\VUgras{大隊} \textsuperscript{(47)} の先頭に立ち、
\VUgras{大尉}はおのおの自分の\VUgras{中隊}を率い、\VUgras{中尉}、\VUgras{少尉}、
\VUgras{下士官}は部下のかたわらの、いつもの位置についていた。

%%K t14-15-1 p
15. --- \VUgras{十字路} \textsuperscript{(50)} には大勢の通行人が群れていた。
商人たち——\VUgras{傘屋} \textsuperscript{(51)}、\VUgras{染物屋} \textsuperscript{(53)}、
\VUgras{鉄砲屋} \textsuperscript{(54)} ——は\VUgras{兵隊}を見に出てきた。車という車も、
\VUgras{辻馬車} \textsuperscript{(52)}、\VUgras{自家用馬車} \textsuperscript{(57)}、それに
\VUgras{散水車} \textsuperscript{(56)} までが、歩道ぞいに寄って停まった。

%%K t14-tit-7 sub
\VUpk{10.2pt}{40}{街路の情景。}
\VUsaut{3.0mm}

%%K t14-15-2 p
\VUgras{行商人} \textsuperscript{(55)} が\VUgras{見本箱}を手にさげて足早に街路を横切り、
\VUgras{屋上席} \textsuperscript{(59)} ——\VUgras{乗合馬車} \textsuperscript{(58)} の屋上席
である——に座っていた人々は見ものを楽しもうと振り返った。気の毒な\VUgras{流しの歌い手}
\textsuperscript{(60)} は\VUgras{手回しオルガン} \textsuperscript{(61)} を連れていたが、
太鼓の音に声を消されて\VUgras{歌}をやめるほかなかった。

%%K t14-16-1 p
16. --- 連隊が\VUgras{羅紗店} \textsuperscript{(64)} の前まで来ると、
\VUgras{お針子} \textsuperscript{(63)} と\VUgras{仕立屋} \textsuperscript{(62)} は
\VUgras{ミシン}と仕事を放り出して窓から眺めた。\VUgras{店主} \textsuperscript{(65)} は
新しい\VUgras{服} \textsuperscript{(66)} を\VUgras{陳列棚}に並べたばかりだったが、
\VUgras{羅紗} \textsuperscript{(67)} と布の巻物——歩道の\VUgras{下水口}
\textsuperscript{(68)} のそばに出してあった——を中に入れさせなければならなかった。
人ごみに倒されるのを恐れたのである。年とった\VUgras{行商人} \textsuperscript{(69)} まで
が、門番の女と靴下の値を掛け合っている最中に、商いを終えるため戸口の廊下に入らねば
ならなかった。

%%K t14-16-2 p
私たちは連隊について兵営まで行き、\VUgras{そして}この一日にたいそう満足して、夕食の
時刻に家へ帰った。

%%K t14-tit-8 sub
\VUpk{10.2pt}{40}{さまざまの商売と職業。}
\VUsaut{3.0mm}

%%K t14-17-1 p
17. --- その次の日、父が土地の新聞の印刷所を見に行こうと言い出した。私たちは
喜んで承知した。

%%K t14-17-2 p
その\VUgras{印刷業者}は\VUgras{書店} \textsuperscript{(70)} \VUgras{（＝本屋）}でもあり、
\VUgras{出版業者}、\VUgras{紙屋}、\VUgras{製本屋}でもある。彼は二階に新聞の
\VUgras{編集室} \textsuperscript{(71)} を置き、\VUgras{製版室}も置いた。そこでは
\VUgras{石版工} \textsuperscript{(72)} と\VUgras{銅版彫刻工} \textsuperscript{(73)} が働いて
いる。最後に\VUgras{植字場}があって、そこには\VUgras{植字工} \textsuperscript{(74)} が
いる。本来の\VUgras{印刷場} \textsuperscript{(75)}、\VUgras{印刷機}、いろいろの機械は
一階にある。

%%K t14-tit-9 sub
\VUpk{10.2pt}{40}{ジャンは質問ばかりする。}
\VUsaut{3.0mm}

%%K t14-18-1 p
18. --- 印刷所を出るとき、ジャンはひどく不思議がって、柱に取りつけた小さな
ガラスの箱の前で足をとめた。\VUgras{小間物屋} \textsuperscript{(77)} の向かいである。彼は
それが何に使うものかと尋ねた。父は\VUgras{火災報知器} \textsuperscript{(76)} だと説明して
やった。じっさい町のものは何もかも私の友だちには驚きだった。素朴な\VUgras{石の水飲み場}
\textsuperscript{(78)} や、軽い\VUgras{配達三輪車} \textsuperscript{(80)} ——お仕着せを着た
\VUgras{給仕} \textsuperscript{(79)} が漕いでいる——から、\VUgras{郵便車}
\textsuperscript{(81)} にいたるまで。

%%K t14-19-1 p
19. --- 父がしばらく私たちを離れて\VUgras{収税吏} \textsuperscript{(83)} と話しこんだ
——その人は\VUgras{弁護士} \textsuperscript{(82)} や\VUgras{公証人} \textsuperscript{(84)}
と立ち話をしていたのである——あいだ、私たちは街路の往来を目で追って退屈をまぎらした。
実にさまざまの人が前を通っていった。\VUgras{菓子屋の小僧} \textsuperscript{(85)} が籠に
みごとな\VUgras{パイ}を入れて、\VUgras{古着屋} \textsuperscript{(86)} のあとから来る。
\VUgras{点灯夫} \textsuperscript{(87)} が\VUgras{ガス工} \textsuperscript{(88)} と話して
いる。\VUgras{修道女} \textsuperscript{(89)} が孤児の女の子の手を引いて修道院へ帰っていく。

%%K t14-tit-10 sub
\VUpk{10.2pt}{40}{家へ帰る道すがら。}
\VUsaut{3.0mm}

%%K t14-20-1 p
20. --- 父が私たちのところへ戻ってきたとたん、ジャンは大笑いした。すすだらけの
\VUgras{煙突掃除人} \textsuperscript{(91)} が二人、汚れた顔と妙ななりをしていたからである。

%%K t14-21-1 p
21. --- 私は\VUgras{花売りの女} \textsuperscript{(92)} から母に鉢植えを買おうと思った
が、父が、持って帰るには重すぎる、\VUgras{オレンジ}にしたほうがよいと言った。そこで
\VUgras{売り子} \textsuperscript{(94)} のところへ寄った。\VUgras{家庭教師}
\textsuperscript{(93)} が教え子のために選んでいたので、その買い物が終わるのを待って
から、私たちも買った。

%%K t14-22-1 p
22. --- 帰り道で知り合いに何人も出会った。家の古い友だちの婦人は\VUgras{連れの婦人}
\textsuperscript{(95)} の腕に手を掛けていた。橋と道路の\VUgras{技師} \textsuperscript{(96)}
にも会った。従姉の一人は\VUgras{子守} \textsuperscript{(98)} を連れていた。

%%K t14-22-2 p
そのあと母の\VUgras{帽子屋} \textsuperscript{(97)} の前を通り、この町の者でない
\VUgras{修道士} \textsuperscript{(99)} 二人に出会った。二人は父に駅へ行く道を尋ねに
やってきたのである。
\VUsaut{6.0mm}
\VUfilet{22mm}
